\documentclass[10pt,a4paper]{article}
\usepackage[utf8]{inputenc}
\usepackage[a4paper,left=3cm,right=2cm]{geometry}
\usepackage{amsmath}
\usepackage[thmmarks, amsmath, thref]{ntheorem}
\usepackage{amsfonts}
\usepackage{amssymb}
\usepackage{fancyhdr}
\pagestyle{fancy}
\usepackage{anysize}
\usepackage{graphicx}
\usepackage{float}
\usepackage{hyperref}
\usepackage{enumitem}
\newtheorem{theorem}{Theorem}
\theoremsymbol{\ensuremath{\blacksquare}}
\newtheorem*{solution}{Solución:}
\usepackage{float}
\usepackage{xcolor}

% Margenes y formalidades

\textwidth 6.25in % Margen Derecho y ancho del texto
\textheight 8.5in
\oddsidemargin 0in % Margen Izquierdo
\headheight 0in % Largo del texto
\leftmargin 1in
\parindent 0em % Salto entre la indentación (Sangría)https://www.overleaf.com/4519127669sdpptkhvsbfn
\topmargin -0.5in % Margen Superior
\parskip 0.5ex % Salto entre parrafos
\baselineskip 1pt



\lhead{
\includegraphics[scale=0.09]{Figuras/logoUSM.png}} % texto izquierda de la cabecera
%\chead{}                                      % texto centro de la cabecera
\rhead{Universidad Técnica Federico Santa María\\Ingeniería Comercial } % número de página a la derecha
\renewcommand{\headrulewidth}{0.4pt} % grosor de la línea de la cabecera
\renewcommand{\footrulewidth}{0.4pt} % grosor de la línea del pie

\begin{document}
\vspace{5cm}





\begin{center}
\begin{large}
\textbf{Econometría ICS-294: Precertamen 3}
\end{large}\\
\vspace{0.2cm}
\end{center}

\section{Preguntas Cortas }

\begin{enumerate}

\item  ¿En qué casos, y por qué motivos, conviene transformar las variables a primeras diferencias cuando se trabaja con series de tiempo?
\\

\item  Defina claramente las dos condiciones que debe cumplir una variable instrumental para ser considerada válida e indique qué problema aparece si se viola cada una de ellas.  \\


\item  Compare el test de White (versión completa) con el test de Breusch-Pagan y mencione una ventaja y una desventaja clave de emplear el test de White.\\

\end{enumerate}

\section{Ejercicios }

\begin{enumerate}

\item Se dispone de una muestra de 350 observaciones con la cual se estimó el siguiente modelo de regresión lineal múltiple:
$$Salario_i = \beta_0 + \beta_1\,\textit{Educ}_i + \beta_2\,\textit{Exper}_i + u_i$$
Se desea contrastar el supuesto de heterocedasticidad y, para ello, se aplicará el método de White simplificado.
\begin{enumerate}
    \item  Plantee la regresión auxiliar para el método de White simplificado.\\
   
    
    \item  Formule la hipótesis nula y la hipótesis alternativa para contrastar la presencia de heterocedasticidad.\\
    
    
    \item  Si el $R^2$ de la regresión auxiliar es $0.14$, contraste, al 5\% de significación, la existencia de heterocedasticidad mediante el estadístico del multiplicador de Lagrange.\\
    
\end{enumerate}

\item  Considere el siguiente modelo de regresión lineal  de los puntajes de una prueba para la selección universitaria.
$$Puntaje_i = \beta_0 + \beta_1 IQ_i + \beta_2 Edad_i + \beta_3 HE$$
Siendo $IQ$ el coeficiente intelectual, la $Edad$ es la edad del estudiante y $HE$ las horas de estudio. Se sospecha que la regresión presenta multicolinealidad en la variable $HE$. 

\begin{enumerate}
    \item Plantee la regresión auxiliar que nos permite encontrar multicolinealidad en la variable $HE$ usando VIF. \\


    \item  Si $R^2_{HE} = 0.91$, calcule VIF y concluya si existe multicolinealidad sobre esta variable.
\end{enumerate}

\item  Se sospecha que la serie de producción de una empresa  sigue un proceso integrado de orden (1). Se plantea la siguiente regresión para testearlo:
$$\Delta y_t = \beta_0 + \phi y_{t-1} + v_t$$
Se encuentran los siguientes resultados para una muestra de 200 datos:
\begin{center}
    \begin{figure}[H]
        \centering
        \includegraphics[width=0.6\linewidth]{certamenes/Resultado Serie de tiempo.jpg}
        \caption{Resultados Regresión}
        \label{fig:enter-label}
    \end{figure}
\end{center}
\begin{enumerate}
    \item  Plantee la hip\'otesis nula y la hipotesis alternativa que permite testear si la serie  tiene ra\'iz unitaria.
    \\
  

    \item  Testee con un 5\% de significancia si la regresión en primeras diferencias tiene raiz unitaria y concluya. \\
  
\end{enumerate}


 \item Un investigador ha simulado tres series temporales de 200 observaciones. Dos de ellas corresponden a procesos \textbf{AR(1)} con distintos valores del parámetro \( \phi \), y la tercera corresponde a un proceso de \textbf{ruido blanco} (white noise). Para cada caso se presentan dos gráficos:

\begin{itemize}
    \item A la izquierda: la serie temporal simulada \( \{X_t\} \)
    \item A la derecha: su función de autocorrelación (ACF) hasta el rezago 30
\end{itemize}

Los valores de \( \phi \) utilizados para los procesos AR(1) fueron:

\[
\phi = -0.95 \quad \text{y} \quad \phi = 0.95
\]

\[
X_t = \phi X_{t-1} + \varepsilon_t
\]

Los gráficos han sido rotulados como \textbf{Caso A}, \textbf{Caso B} y \textbf{Caso C} (sin indicar el valor de \( \phi \) ni el tipo de proceso que representan).
\begin{figure}[H]
    \centering
    \includegraphics[width=0.7\linewidth]{certamenes/autocorr.png}
\end{figure}

\begin{enumerate}
    \item  Identifique cuál de los tres casos corresponde al proceso de \textbf{ruido blanco}. Justifique su respuesta con base en el comportamiento de la serie y su ACF.


% \textcolor{blue}{
%         El \textbf{Caso B} corresponde al proceso de ruido blanco. La serie temporal no muestra una estructura clara ni dependencia entre observaciones sucesivas. Además, (ANALISIS MÁS IMPORTANTE PUEDE TENER TODO EL PUNTAJE CON ESTA PARTE) la función de autocorrelación (ACF) muestra valores cercanos a cero para todos los rezagos, lo cual es característico de un proceso de ruido blanco.
%         }

    \item Para los otros dos casos, asigne a cada uno el valor de \( \phi \) correspondiente. Justifique su elección explicando cómo el signo y la magnitud de \( \phi \) se reflejan en la dinámica de la serie y en la forma de su ACF.

% {\color{blue}
%         El \textbf{Caso A} corresponde al proceso AR(1) con \( \phi = -0.95 \). Se observa un patrón oscilante en la serie, donde los valores tienden a alternar de signo. Esta dinámica es típica de un proceso AR con un coeficiente negativo cercano a -1. La ACF también muestra un patrón alternado (signo negativo en rezagos impares) (ANALISIS MÁS IMPORTANTE PUEDE TENER TODO EL PUNTAJE CON ESTA PARTE) , decreciendo gradualmente, lo que confirma el signo negativo de \( \phi \).

%         El \textbf{Caso C} corresponde al proceso AR(1) con \( \phi = 0.95 \). La serie muestra una fuerte persistencia: los valores permanecen en la misma dirección por varios periodos antes de cambiar, lo cual indica una alta autocorrelación positiva. La ACF (ANALISIS MÁS IMPORTANTE PUEDE TENER TODO EL PUNTAJE CON ESTA PARTE)  decrece lentamente y de forma positiva, lo que es consistente con un valor de \( \phi \) positivo cercano a 1.
%         }   
\end{enumerate}



\end{enumerate}





\section*{Formulario}
\begin{itemize}

\item Tabla $\chi^2$
\begin{center}
    \begin{tabular}{|c|c|c|c|}
    \hline
        df & 0.99 & 0.95 & 0.90 \\ 
        \hline
        1 &  6.635 & 3.841 & 2.706\\
        \hline
        2 &  9.210 & 5.991 & 4.605 \\
        \hline
        3 & 11.345 & 7.815 & 6.251  \\
        \hline
        4 &  13.277 & 9.488 & 7.779 \\
        \hline
        5 &  15.086 & 11.070 & 9.236 \\
        \hline
        6 &  16.812 & 12.592 & 10645 \\
        \hline
    \end{tabular}
\end{center}

\item Multiplicado de Lagrange

$$\color{red}{LM = nR^2 \sim \chi^2}$$

\item Test F para regresión auxiliar.
\begin{center}
 {  \color{red}     \(F=\dfrac{R_{\text{aux}}^{2}/q}{(1-R_{\text{aux}}^{2})/(n-k-1)}
          \sim F_{q,\;n-k-1}\).}
\end{center}



\item Tabla de test D-F
\begin{center}
    \begin{figure}[H]
        \centering
        \includegraphics[width=0.7\linewidth]{sem14/tabla.png}
    \end{figure}
\end{center}

\end{itemize}


\end{document}